\subsubsection*{Численный пример}

\begin{idea}[title={Откуда этот материал}]
Прогон метода отсечений на числах; на лекции пример не разбирался. Задача
специально взята в том виде, к которому алгоритм применим без преобразований:
цель линейна, ограничение одно и выпукло.
\end{idea}

\begin{example}\label{ex:cutting-plane-run}
\[
  \begin{cases}
    \varphi_0(x) = x_1 \to \min , \\
    \varphi(x) = x_1^2 + x_2^2 - 1 \le 0 ,
  \end{cases}
  \qquad
  \Omega = \{\, x \mid \norm{x} \le 1 \,\} ,
  \qquad
  c = \begin{pmatrix} 1 \\ 0 \end{pmatrix} .
\]
Ответ очевиден: $x_* = (-1, 0)\T$, $\varphi_0(x_*) = -1$. Начальный объемлющий
многогранник — квадрат
\[
  S_0 = \{\, x \mid -2 \le x_1 \le 2,\ -2 \le x_2 \le 2 \,\} \supset \Omega .
\]
Градиент ограничения $\nabla\varphi(x) = (2x_1,\ 2x_2)\T$ — он же субградиент,
так как $\varphi$ дифференцируема.

\medskip
\noindent\textbf{Итерация 0.} Линейная подзадача $x_1 \to \min$, $x \in S_0$
даёт $x_1 = -2$ при любом $x_2 \in [-2,2]$: минимум достигается на всём левом
ребре. Выберем $x_0 = (-2, 0)\T$ (см. замечание после примера). Проверка
допустимости:
\[
  \varphi(x_0) = 4 + 0 - 1 = 3 > 0 ,
\]
точка недопустима, строим отсечение:
\[
  a_1 = \nabla\varphi(x_0) = \begin{pmatrix} -4 \\ 0 \end{pmatrix} ,
  \qquad
  b^{(1)} = -\varphi(x_0) + a_1\T x_0 = -3 + 8 = 5 ,
\]
\[
  a_1\T x - b^{(1)} \le 0
  \quad\Longleftrightarrow\quad
  -4x_1 - 5 \le 0
  \quad\Longleftrightarrow\quad
  x_1 \ge -\tfrac54 .
\]
\[
  S_1 = S_0 \cap \big\{\, x_1 \ge -\tfrac54 \,\big\} .
\]
Отсечение корректно: для любой точки круга $\abs{x_1} \le 1$, то есть
$x_1 \ge -1 > -\tfrac54$, — ни одна допустимая точка не потеряна
(лемма~\ref{lem:cut-nested}), а сама $x_0$ выброшена.

\medskip
\noindent\textbf{Итерация 1.} Решение подзадачи на $S_1$:
$x_1 = (-\tfrac54,\ 0)\T$,
\[
  \varphi(x_1) = \tfrac{25}{16} - 1 = \tfrac{9}{16} > 0 ,
\]
\[
  a_2 = \nabla\varphi(x_1) = \begin{pmatrix} -\tfrac52 \\ 0 \end{pmatrix} ,
  \qquad
  b^{(2)} = -\tfrac{9}{16} + \Big(-\tfrac52\Big)\Big(-\tfrac54\Big)
          = -\tfrac{9}{16} + \tfrac{50}{16} = \tfrac{41}{16} ,
\]
\[
  -\tfrac52 x_1 - \tfrac{41}{16} \le 0
  \quad\Longleftrightarrow\quad
  x_1 \ge -\tfrac{41}{40} = -1{,}025 ,
  \qquad
  S_2 = S_1 \cap \big\{\, x_1 \ge -\tfrac{41}{40} \,\big\} .
\]

\medskip
\noindent\textbf{Итерация 2.} $x_2 = (-\tfrac{41}{40},\ 0)\T$,
\[
  \varphi(x_2) = \tfrac{1681}{1600} - 1 = \tfrac{81}{1600} \approx 0{,}0506 .
\]
Невязка ограничения падает быстро:
\[
  3 \ \to\ \tfrac{9}{16} = 0{,}5625 \ \to\ \tfrac{81}{1600} \approx 0{,}0506 ,
\]
а приближения к оптимуму идут снизу:
\[
  -2 \ \to\ -1{,}25 \ \to\ -1{,}025 \ \to\ \dots \ \to\ -1 .
\]
Все $x_k$ недопустимы, поэтому $\varphi_0(x_k) \le \varphi_0(x_*)$ — внешняя
аппроксимация даёт \emph{нижние} оценки оптимума, что и использовалось в
доказательстве теоремы~\ref{thm:cutting-plane}.
\end{example}

\begin{remark}[о неединственности подзадачи]
Минимум $x_1$ на $S_k$ достигается на целом отрезке, и симплекс-метод вернул бы
какую-то вершину, например $(-2,-2)\T$. Мы фиксируем $x_2 = 0$: задача
симметрична относительно оси $x_1$, и такой выбор удерживает все итерации на
одной прямой, что делает счёт обозримым. На скорость это не влияет — отсечение
всё равно получается вертикальной прямой $x_1 = \text{const}$.
\end{remark}

\begin{idea}
Любопытная деталь: на прямой $x_2 = 0$ ограничение превращается в
$g(t) = t^2 - 1 \le 0$, а отсечение — в
\[
  g(t_k) + g'(t_k)(t - t_k) \le 0
  \quad\Longleftrightarrow\quad
  t \ \ge\ \frac{t_k^2 + 1}{2 t_k}
  \qquad (t_k < 0) ,
\]
то есть $t_{k+1} = \tfrac12\big(t_k + \tfrac{1}{t_k}\big)$ — в точности метод
Ньютона для $\sqrt{1}$. Отсюда и наблюдаемая скорость: $-2$, $-1{,}25$,
$-1{,}025$, $-1{,}0003$ — число верных знаков удваивается за шаг. В общем
случае такой скорости у метода отсечений нет: она здесь берётся из гладкости и
одномерности, а не из самого метода.
\end{idea}
